\paragraph{AQ-136-RELATIVE-H1: a relative integral repair through vertical degree one.}
This is a positive partial construction in the original strict PD
category. It does not assert the higher vertical vanishings required
for a fully admissible source.

Let $k$ be a nonzero commutative ring which is torsion-free as a
$\mathbb Z$-module, and put $K=k\otimes_{\mathbb Z}\mathbb Q$.
Let $C_1,\ldots,C_t\in\mathcal A_k$, where $t\ge1$.
Suppose $A$ is a simplicial nonnegatively bigraded strict PD
differential $S_k$-algebra, with maps $a_j:A\to C_j$, such that
\[
 A_0=S_k,\qquad A_n^{0,*}=k[0],\qquad
 H^0((A_K)_n^{w,*})=S_K^w,\quad
 H^1((A_K)_n^{w,*})=0\quad(w\ge W).                       \tag{136.RH1}
\]
The coefficient identification in (136.RH1) is canonical and
holds uniformly in $n$. Enlarge $W\ge1$ so that for every target
and every $n$ one also has
\[
 H^0((C_j)_n^{w,*})=S_k^w,
       \qquad H^1((C_j)_n^{w,*})=0\quad(w\ge W).           \tag{136.RH2}
\]
There is a sequence of strict maps under $A$ and over all targets
whose filtered colimit $A^+$ satisfies
\[
 \begin{gathered}
 A^+_0=S_k,\qquad (A^+)_K\simeq A_K
                   \quad\text{as actual differential PD algebras},\\
 H^0((A^+)_n^{w,*})=S_k^w,\qquad
 H^1((A^+)_n^{w,*})=0\quad(w\ge W,\ n\ge0).
 \end{gathered}                                             \tag{136.RH3}
\]
No horizontal acyclicity is claimed for this repair by itself.
The same single bound $W$ works throughout and in the colimit.
The linear inputs are the Dold--Kan equivalence
(Stacks Project, Theorem~14.24.3, Tag~019G) and the free PD
polynomial universal property (Lemma~23.5.1, Tag~07H5).
The arbitrary-module extension of the latter is constructed
explicitly below; no flatness hypothesis is imported from it.

\begin{proof}
We use normalization $N_r$ in the simplicial direction, with
horizontal differential $\partial:N_r\to N_{r-1}$, and retain
$d$ for the form differential. Normalization is exact and its
individual terms commute with vertical cohomology: each is a
natural direct summand of the corresponding simplicial term.
Thus (136.RH2) gives, for every $r\ge1$ and $w\ge W$,
\[
 H^0(N_r C_j^{w,*})=H^1(N_r C_j^{w,*})=0.
                                                               \tag{136.RH4}
\]
In particular every vertical one-cycle in $N_rC_j$ has a
\emph{unique} zero-form primitive. The corresponding rational
statement holds for $N_r A_K$ by (136.RH1).

We first define one repair step $A\to R(A)$.
In each weight $w\ge W$ define a nonnegative horizontal chain
complex $V^w$ by
\[
 V_0^w=S_k^w,\qquad
 V_r^w=Z^1(N_r A^{w,*})\quad(r\ge1).                       \tag{136.RH5}
\]
For $r\ge2$ its differential is the existing $\partial$ on
one-cycles. It remains to define $\partial:V_1^w\to V_0^w$.
If $u\in Z^1(N_1 A^{w,*})$, rational exactness and the absence
of normalized degree-one vertical $H^0$ give a unique
\[
 v_K\in N_1 A_K^{w,0},\qquad d v_K=u_K.
\]
The image $a_j(u)$ has, by (136.RH4), a unique integral primitive
$y_j\in N_1 C_j^{w,0}$. After scalar extension its value is
$a_{j,K}(v_K)$, by rational uniqueness. Since $A_0^{w,1}=0$,
\[
             d(\partial y_j)=a_j(\partial u)=0.
\]
Consequently $\partial y_j$ is the image of a unique element
$s_j\in S_k^w$, by (136.RH2) and nonnegative form degrees.
Its rationalization is the element $\partial v_K$ of
$A_{K,0}^{w,0}=S_K^w$. The injection $S_k^w\hookrightarrow S_K^w$
therefore makes all $s_j$ equal. Define
\[
                         \partial u=s_j\quad\text{in }V_0^w.
                                                               \tag{136.RH6}
\]
It is $k$-linear: rational primitives and their boundaries are
linear, and the coefficient injection is injective.
For $u\in V_2^w$, the rational primitive of $\partial u$ is
$\partial$ of the primitive of $u$. Its next boundary is zero.
The same injection proves $\partial^2u=0$ integrally.
Thus (136.RH5)--(136.RH6) is a chain complex.

For $w<W$ instead put $V^w=N(A^{w,0})$ unchanged.
There are horizontal chain maps
\[
 f:N(A^{*,0})\longrightarrow V,
             \qquad d_V:V\longrightarrow N(A^{*,1}),      \tag{136.RH7}
\]
where in high weights $f_0$ is the coefficient identity,
$f_r(a)=da$ for $r\ge1$, $(d_V)_0=0$, and $(d_V)_r(u)=u$
for $r\ge1$. In low weights use the identity and the old
vertical differential. The only boundary identity requiring
comment is $r=1$: the rational primitive of $da$ is $a_K$,
so (136.RH6) gives precisely $\partial a\in S_k^w$.
Thus
\[
                         d_Vf=d_A,\qquad d_A d_V=0.
                                                               \tag{136.RH8}
\]

Each target map on zero forms extends to a chain map
$g_j:V\to N(C_j^{*,0})$. In high normalized degrees $r\ge1$,
send $u$ to its unique primitive from (136.RH4); in normalized
degree zero use coefficients. For $r\ge2$ compatibility with
$\partial$ follows from uniqueness of the primitive in degree
$r-1$, and for $r=1$ it is the definition (136.RH6).
The old image of $a\in N_rA^{w,0}$ is itself the primitive
of $a_j(da)$, so
\[
                 g_j f=N(a_j^0),\qquad
                   d_{C_j}g_j=N(a_j^1)d_V.                \tag{136.RH9}
\]
Dold--Kan now gives actual simplicial modules and maps with
all these identities. We continue to denote them by $V,f,d_V,g_j$.

In positive weights take the graded PD algebra pushout
\[
 R(A)=A\amalg_{\Gamma_k(A^{*,0}_{>0})}
                              \Gamma_k(V_{>0}).           \tag{136.RH10}
\]
This formula first denotes a pushout of graded algebras;
the middle algebra is not being given a zero differential
and mapped to $A$ as a differential algebra. Instead define
the differential on the new zero-form generators by
$d[v]=d_V(v)$, and keep the old differential on $A$.
The identities (136.RH8) make every identifying relation
$[f(a)]-a$ stable under this differential. Its PD ideal is
also stable by the special derivation rule. The differential
therefore descends to (136.RH10). Its square vanishes on
old elements and new generators; on divided powers it also
vanishes, using the square-zero rule for the odd element
$d[v]$. This gives a strict PD differential algebra with
the original nonnegative form grading. All constructions
are simplicial and preserve weights and the coefficient map.
The maps (136.RH9) and the free PD universal property give
strict maps $R(A)\to C_j$ extending all $a_j$.
The construction can equivalently be made as the free graded
PD extension by $V_{>0}$ followed by its displayed identifying
PD ideal, so it requires no flatness of $V$ or of a target module.
For an arbitrary module the free PD algebra itself is obtained
from a PD polynomial algebra on a set of module generators by
imposing the PD ideal of all linear relations. This also proves
its universal property and compatibility with scalar extension.

In simplicial degree zero $V_0=S_k=A_0$, so (136.RH10)
changes nothing there. Weight zero also remains $k[0]$.
After extension to $K$, the map $f$ is an isomorphism.
For high weights and $r\ge1$, this is exactly the bijection
\[
 d:N_rA_K^{w,0}\xrightarrow{\sim} Z^1(N_rA_K^{w,*})
\]
from (136.RH1); localization commutes with cycles.
In other weights and normalized degree zero it is already
an identity. Dold--Kan and the free PD universal property,
which commutes with scalar extension, now identify
\[
                           R(A)_K\simeq A_K               \tag{136.RH11}
\]
as differential PD algebras, compatibly with the transition.

The step kills every previous vertical one-cycle in high
weights: in positive normalized degree $r$ the new generator
$[u]$ has differential $u$. In normalized degree zero there
are no one-forms because $A_0=S_k$. It also kills every
previous normalized vertical zero-cycle in degree $r\ge1$,
since its image under $f$ is $da=0$. Consequently the image
of $H^0(A_n^{w,*})\to H^0(R(A)_n^{w,*})$ consists of
coefficients, and the map on $H^1$ is zero, for $w\ge W$.
The normalized direct-sum decomposition and its degeneracies
justify these statements in every simplicial degree, not only
on normalized terms.

Iterate $R$ countably many times and take the filtered colimit.
Every step still satisfies (136.RH1), with the same $W$, by
(136.RH11), and retains $A_0=S_k$ and the target maps.
Filtered colimits of modules are exact and commute with the
finite kernel/cokernel calculations for fixed vertical
cohomology. The same is true of the underlying modules of
these algebras: each operation and each defining identity
has finite arity, including each fixed divided power.
The transition kills every $H^1$ class at the next step and
sends every $H^0$ class to a coefficient. Coefficients remain
injective, since simplicial degeneracies from degree zero have
face retractions and degree zero is still exactly $S_k$.
This proves all assertions of (136.RH3), including the single
uniform bound in all simplices. The rationalized transitions
are isomorphisms compatible with $A_K$, so their colimit is
that actual algebra, as stated.
\end{proof}

\paragraph{AQ-136-PARTIAL-SOURCE: all horizontal rows and the first two vertical conditions.}
For the same torsion-free $k$ and any finite nonempty family
$C_1,\ldots,C_t\in\mathcal A_k$, there is a simplicial strict
PD differential $S_k$-algebra $A$ with strict maps to all $C_j$
such that:
\[
 \begin{gathered}
 A_0=S_k,\qquad A_n^{0,*}=k[0],\qquad H^*(A_{K,n})=S_K
                                   \text{ in every weight},\\
 H^0(A_n^{w,*})=S_k^w,\quad H^1(A_n^{w,*})=0
                        \quad(w\ge W,\ n\ge0),\\
 H_r(N(A_*^{w,q}))=0
                   \quad(w\ge\mu(q),\ r\ge0,\ q\ge0).
 \end{gathered}                                             \tag{136.RH12}
\]
The bounds depend on the given targets, not on the construction
stage or simplex. All positive vertical cohomology groups are
integer-torsion; no vanishing in vertical degree $q\ge2$ is
asserted. In particular (136.RH12) is not yet an object of
$\mathcal A_k$ merely from the conclusions proved here.

\begin{proof}
Choose $W$ from (136.RH2) for the targets. For every form degree
$q$ choose $\mu(q)\ge1$ so that all their normalized horizontal
rows are acyclic at weights $w\ge\mu(q)$.
Begin with the constant coefficient algebra $A^{(0)}=S_k$
and its given maps to all targets. We describe a horizontal
cell step $A\to H(A)$ preserving $A_0=S_k$, weight zero,
and the rational vertical comparison in every weight.

Let $z\in N_r A^{w,q}$ be a horizontal cycle, where
$w\ge\mu(q)$. A free vertical disk $E_q$ has generators
$a,b$ of form degrees $q,q+1$ and $da=b$, $db=0$.
Give them both weight $w$. Let $\mathbb S^r$ be the
simplicial module corresponding to $k[r]$ and let
$\mathbb D^{r+1}$ correspond to the horizontal disk
$k[r+1]\xrightarrow{1}k[r]$. The boundary inclusion
$\mathbb S^r\hookrightarrow\mathbb D^{r+1}$ is split in
each simplicial module degree. The pair $z,dz$ specifies a
map from the free graded PD differential algebra on
$\mathbb S^r\otimes E_q$ to $A$. Attach the free algebra on
$\mathbb D^{r+1}\otimes E_q$ by a pushout along this map.
In normalization the new generator $a$ has
$\partial a=z$, and its partner satisfies
$\partial b=dz$.

For each target the image of $z$ is a boundary in its form-$q$
row. Choose a normalized primitive $y_j$ with
$\partial y_j=a_j(z)$, and send the new generators to
$y_j,dy_j$. This extends the target maps as strict simplicial
PD differential maps. The free universal property checks all
simplicial identities simultaneously. Attach such cells for
all cycles and all triples $(r,w,q)$ satisfying the bound.
These cells form a set, and a simultaneous coproduct and
pushout therefore exist.

No new generator occurs in simplicial degree zero: for $r=0$
the boundary inclusion is already an isomorphism there, and
for $r>0$ both degree-zero modules vanish. Thus $A_0=S_k$
remains exact. Positive weights of the cells retain weight zero.
At a fixed simplex the boundary-module inclusion splits as a
module, and the splitting tensored with $E_q$ respects its
vertical differential. The attached algebra is consequently
an extension of $A_n$ by free vertical PD disks, with their
new generator pairs independent over $A_n$.
After extension to $K$, each such disk is contractible onto
its constants. One explicit proof uses the degree-minus-one
special derivation sending $b$ to $a$ and $a$ to zero:
its commutator with $d$ is multiplication by the positive
monomial arity. Divide by that arity in the $\mathbb Q$-algebra
$K$ to obtain a contraction on the augmentation ideal.
The contraction preserves weight. Tensoring with $A_n$
and taking filtered colimits over finite subsets of the new
generators proves that the attachment preserves rational
vertical cohomology in every weight and simplex.
By construction the transition induces zero on every old
horizontal homology group in the specified tails.

Now alternate one horizontal step $H$ with one repair step
$R$ of the preceding theorem, always using the same $\mu(q)$
and $W$, and take the countable filtered colimit. Every stage
still has simplex zero $S_k$ and exact rational vertical
cohomology, so the repair is always applicable. Every old
horizontal class in the weight tail is killed at the next
horizontal step and remains zero afterwards. Every old
vertical $H^1$ class is killed at the next repair, and every
old vertical $H^0$ class maps to a coefficient there.
Exactness of filtered colimits gives all the assertions of
(136.RH12), with the original fixed bounds.
Finally scalar extension to $K$ commutes with cohomology and
has zero positive vertical cohomology. Its kernel on any
underlying abelian group consists exactly of elements killed
by some nonzero integer. Thus all remaining positive vertical
cohomology is integer-torsion. No common annihilator or higher
vanishing has been deduced.
\end{proof}
